\documentclass[conference]{IEEEtran}
\IEEEoverridecommandlockouts

\usepackage{cite}
\usepackage{amsmath,amssymb}
\usepackage{graphicx}
\usepackage{booktabs}
\usepackage{xcolor}
\usepackage{url}

% A figure from figures/ if it is there, else a framed placeholder, so the draft
% compiles before the figures from paper_results.zip are copied in.
\newcommand{\fig}[2]{%
  \IfFileExists{figures/#2}{\includegraphics[width=#1]{figures/#2}}%
  {\fbox{\parbox[c][3cm][c]{0.9#1}{\centering\small figures/\detokenize{#2}\par (copy in from paper\_results.zip)}}}}

\begin{document}

\title{One Predictor for Four Rail Faults from Vibration and Sound Alone}

\author{
\IEEEauthorblockN{Amarsh N Choudhary}
\IEEEauthorblockA{\textit{Department of Mechanical Engineering} \\
\textit{Indian Institute of Technology Delhi}\\
New Delhi, India \\
me2252094@iitd.ac.in}
\and
\IEEEauthorblockN{Mayank Tayal}
\IEEEauthorblockA{\textit{Department of Electrical Engineering} \\
\textit{Indian Institute of Technology Delhi}\\
New Delhi, India}
\and
\IEEEauthorblockN{Himanshu}
\IEEEauthorblockA{\textit{Department of Electrical Engineering} \\
\textit{Indian Institute of Technology Delhi}\\
New Delhi, India}
}

\maketitle

\begin{abstract}
Cameras and track-recording cars are the usual tools for finding rail defects, but sound and
vibration can be captured with cheap sensors on any train. We build a single fault predictor
for four rail-fault tasks using only microphone and accelerometer data from three public
datasets: rail-surface faults (superelevation and wheel burns), rail corrugation heard by an
in-cabin microphone, corrugation felt at the axle box, and rail joints and switches. For each
task we compare a random forest on spectral features with a fine-tuned Audio Spectrogram
Transformer (AST), and the predictor uses whichever is better. Clips are cut on a distance
grid rather than in time, so a fault produces the same pattern at any train speed. Over
repeated training runs the predictor reaches a test macro-F1 of $0.99$ on rail-surface faults,
$0.92$ for corrugation from the axle box and $0.89$ from the microphone, both measured at a train
speed not seen in training, and $0.81$ on joints and switches. Labels for the joint task come
from an unsupervised detector that places all nine surveyed joints and switches of the
Rail-VIVID dataset inside repeatable stretches of high vibration. We also report a negative
result: a single shared AST with one output head per fault was unstable across four training
runs, and no run matched the separate models on more than one task.
\end{abstract}

\begin{IEEEkeywords}
railway track monitoring, fault detection, vibration, acoustic analysis, audio spectrogram
transformer, rail corrugation, multi-task learning
\end{IEEEkeywords}

%% ------------------------------------------------------------------------------------------
\section{Introduction}

Rail defects such as corrugation, wheel burns, worn joints and poorly set superelevation raise
noise, wear and maintenance cost, and some lead to derailment if they go unnoticed. Most
inspection relies on dedicated measurement trains or on cameras, which are costly to run often.
Microphones and accelerometers are cheap and can ride on ordinary service trains, so the signal
they record could flag faults between scheduled inspections.

Public labelled data for this is scarce and fragmented. Each available dataset covers one line,
one vehicle and one or two kinds of fault, recorded with different sensors at different
sampling rates. This paper asks how far one can get with only that public data and only
vibration and sound, with no images.

Our contributions are:
\begin{itemize}
  \item One predictor that scores four rail-fault tasks from vibration and sound, assembled from
        the best model per task, with results averaged over repeated training runs on grouped
        splits that keep every recording (or run) in a single split.
  \item Distance-normalised clips: each recording is resampled onto a distance grid at a fixed
        reference speed, so a periodic fault such as corrugation keeps the same frequency at any
        train speed. With this, corrugation is detected on a run at a speed not seen in
        training.
  \item An unsupervised way to label joints and switches in Rail-VIVID, whose vibration
        recordings carry no labels: repeatable stretches of high vibration contain all nine
        surveyed joints and switches.
  \item A negative result: a shared multi-task AST was unstable across runs, and no run matched
        the separate models on more than one task.
\end{itemize}
Code to rebuild every result is public (Section~\ref{sec:availability}).

%% ------------------------------------------------------------------------------------------
\section{Related Work}

\textit{Acoustic track-fault detection.} Shafique et al.~\cite{shafique2021} recorded track sound
with microphones on a cart on Pakistan Railways and classified normal track, superelevation and
wheel burns with MFCC features and classical classifiers. We use their data, but evaluate on a
split grouped by recording so that windows of one recording never sit on both sides of the split.

\textit{Corrugation.} Rail corrugation is a periodic wear pattern on the rail head whose
wavelength is set by the track and vehicle, so its frequency in the measured signal rises with
train speed~\cite{grassie1993}. Soto-Ocampo et al.~\cite{soto2025} detect corrugation from
on-board acoustic measurements with a mathematical method and publish the recordings with the
corrugated sections they confirmed by inspection~\cite{upmcorrugation}. We use these recordings
and the reported sections as labels.

\textit{Rail-VIVID.} Ma et al.~\cite{ma2026railvivid} released synchronised accelerometer, camera
and GNSS data from twenty runs over the same 1.4\,km of track, with nine surveyed anomalies
(rail joints and switch points) given by position. We use only the vibration data.

\textit{Pretrained audio models.} The Audio Spectrogram Transformer (AST)~\cite{gong2021ast},
pretrained on AudioSet~\cite{gemmeke2017audioset}, and CNNs such as PANNs~\cite{kong2020panns} are
strong starting points for small sound-classification datasets. We fine-tune AST.

\textit{Multi-task learning.} Sharing one encoder across tasks can help or hurt depending on
which tasks are trained together~\cite{standley2020}. Our shared model is an example of the
latter.

%% ------------------------------------------------------------------------------------------
\section{Data}
\label{sec:data}

Table~\ref{tab:data} summarises the three datasets and the four tasks drawn from them.

\begin{table*}[t]
\centering
\caption{Datasets and tasks. Every split is grouped, so all clips from one recording (or one
run) stay in the same split. Clip counts are after cutting.}
\label{tab:data}
\begin{tabular}{llllrrl}
\toprule
Task & Dataset & Sensor & Classes & Clips & Test & Split grouped by \\
\midrule
Rail surface       & Pakistan Railways~\cite{shafique2021} & microphone, 16\,kHz & normal / superelev. / wheel burn & 720 & 119 & recording block \\
Corrugation (mic)  & UPM corrugation~\cite{upmcorrugation} & cabin microphone, 16\,kHz & normal / corrugated & 694 & 174 & run (speed) \\
Corrugation (axle) & UPM corrugation~\cite{upmcorrugation} & axle-box accel., 4\,kHz & normal / corrugated & 694 & 174 & run (speed) \\
Joint / switch     & Rail-VIVID~\cite{ma2026railvivid}     & accel., 2\,kHz & normal / joint & 2831 & 658 & repeat \\
\bottomrule
\end{tabular}
\end{table*}

\textit{Rail surface.} 720 clips of 17\,s (240 per class) recorded at 22.05\,kHz and resampled to
16\,kHz. Recordings whose file numbers fall in the same block of ten are grouped together,
since neighbouring recordings may come from the same stretch of track; groups are split
70/15/15 within each class (481/120/119 clips).

\textit{Corrugation.} Four runs of an instrumented train over the same 2.4\,km at 40, 50, 60 and
75\,km/h. The recordings carry kilometre positions but no labels; we label as \emph{corrugated}
any clip lying at least 75\% inside one of the two sections reported as severely corrugated
in~\cite{soto2025} (km 4.855--4.914 and 5.239--5.326), and as \emph{normal} any clip lying at
least 20\,m from both. Clips in between are dropped. The 40 and 50\,km/h runs are used for
training, 60\,km/h for validation and 75\,km/h for testing, so every test score is measured at a
speed the model never saw. The microphone and axle-box signals are kept as separate tasks.
Since the corrugated track is only about 146\,m long, corrugated clips overlap (2\,m hop); the
test run has 62 corrugated and 112 normal clips.

\textit{Joints and switches.} Rail-VIVID provides 20 runs (five speed and direction settings,
four repeats each) of six accelerometers at 2\,kHz with GNSS. Section~\ref{sec:vivid} explains
how we label it. We keep the two channels whose vibration repeats between runs and export them
as separate clips of the same group. Repeats 1--2 are used for training, 3 for validation and 4
for testing (470 joint and 188 normal test clips).

%% ------------------------------------------------------------------------------------------
\section{Method}

\subsection{Distance-normalised clips}
A corrugation wavelength $\lambda$ produces a tone at $f = v/\lambda$, so the same defect sounds
different at each train speed $v$. Before cutting clips, we resample every run onto a uniform
distance grid using its position record, as if the train had always run at a reference speed
$v_\mathrm{ref}$. A clip is then a fixed length of track (20\,m), and a defect produces the same
frequency at any speed. We use $v_\mathrm{ref}=20$\,m/s for the UPM data (a 20\,m clip lasts 1\,s)
and $6.5$\,m/s for Rail-VIVID, whose runs are between 5 and 8.4\,m/s. An earlier version that cut
clips in time predicted every 75\,km/h test clip as normal; with distance normalisation the same
models detect it.

\subsection{Labelling joints and switches without labels}
\label{sec:vivid}
The Rail-VIVID vibration files have no per-sample labels; the nine anomalies are given as GNSS
positions. We project every sample onto the arc length along a reference trace of the track and
compute, per metre and per run, band powers and slow-signal features for each channel. Joints
did not appear as isolated impacts on these sensors. What did repeat in every run were stretches
of track, typically 15--30\,m long, with distinctly higher 500--990\,Hz vibration on two channels.
Along the track we found 14 such stretches; each of the nine surveyed anomalies lies inside one
of them, and the other five stretches hold no listed anomaly. Those five may be features of the
track that the survey does not list.

To find the stretches without using the labels, we keep only features whose per-metre profile
correlates between runs. Three features pass this test: 500--990\,Hz band power on two channels and
the slow ($<2$\,Hz) signal on a third. A step detector then marks places where the median of these
features over the next 6\,m differs from the previous 6\,m by more than four noise units. It
returns 30 candidate sites. Every one of the nine surveyed anomalies has a site within 15\,m, and
15 sites lie away from any listed anomaly. Since the survey lists only nine points on 1.4\,km of
track, we read this as a high-recall detector of track features rather than a precise one.

We then label a 20\,m clip as \emph{joint} when its centre is within 6\,m of a surveyed anomaly,
and as \emph{normal} when it is at least 25\,m from every anomaly and from every detected
stretch that holds none (these may be unlisted joints, so they are left out rather than guessed).
Stretches where a channel is stuck at a constant value are skipped. This gives 1995 joint and 836
normal clips.

\subsection{Models}
\textit{Random forest.} Each window is described by 137 features: the mean and standard
deviation of 20 MFCCs~\cite{davis1980} and their deltas, statistics of RMS level, spectral
centroid, bandwidth and roll-off, and 48 log-spaced Welch~\cite{welch1967} band powers relative
to their mean plus the overall level. The band powers keep narrow spectral peaks such as a
corrugation tone, which MFCC means tend to smear. A forest of 400 trees with balanced class
weights~\cite{breiman2001} is trained on 1\,s windows with 0.5\,s hop.

\textit{AST.} We fine-tune AST pretrained on AudioSet~\cite{gong2021ast,gemmeke2017audioset}. Its
position embeddings are cut to the window length (5\,s for rail surface, 1\,s for microphone
corrugation) rather than padding to 10\,s. Each epoch draws eight random windows per recording
with class-balanced sampling, plus speed perturbation ($\pm 15\%$), random gain, additive noise at
10--30\,dB SNR, and SpecAugment masks~\cite{park2019specaugment}. Training uses AdamW with learning
rates of $2\times10^{-5}$ for the encoder and $10^{-3}$ for the head, a cosine schedule with 10\%
warm-up, label smoothing 0.1, at most 15 epochs, and early stopping after 4 epochs without
validation gain. The checkpoint with the best validation macro-F1 is kept.

\subsection{Clip scores and decision threshold}
A clip's class probabilities are the mean over its windows. For the two-class tasks, the
decision threshold on the fault probability is chosen to maximise macro-F1 on the validation
split; with few fault clips, the default 0.5 threshold can miss the fault class entirely. All
reported scores are clip-level test macro-F1 with that tuned threshold, or argmax for the
three-class rail-surface task.

\subsection{The predictor}
The predictor loads one model per task, the AST or forest that scored best, and returns a
decision and class probabilities for each task the input sensor supports: microphone clips are
scored for rail-surface faults and microphone corrugation, accelerometer clips for axle-box
corrugation and joints. Models are trained on separate datasets, so each task only compares a
fault with normal track recorded the same way.

%% ------------------------------------------------------------------------------------------
\section{Results}

Every model was trained with different random seeds (three seeds for the forests, two for the
AST models) on the same fixed splits. Table~\ref{tab:main} gives the mean and standard deviation
of the test macro-F1 and Fig.~\ref{fig:scores} shows the spread.

\begin{table}[t]
\centering
\caption{Test macro-F1, mean $\pm$ standard deviation over seeds (3 for random forests, 2 for
AST). The last column is the model the predictor uses.}
\label{tab:main}
\begin{tabular}{lccc}
\toprule
Task & Forest & AST & Predictor uses \\
\midrule
Rail surface       & 0.958 $\pm$ 0.009 & \textbf{0.992 $\pm$ 0.012} & AST \\
Corrugation (axle) & \textbf{0.924 $\pm$ 0.009} & -- & Forest \\
Corrugation (mic)  & 0.869 $\pm$ 0.028 & \textbf{0.893 $\pm$ 0.038} & AST \\
Joint / switch     & \textbf{0.809 $\pm$ 0.020} & -- & Forest \\
\bottomrule
\end{tabular}
\end{table}

\begin{figure}[t]
\centering
\fig{\columnwidth}{scores.pdf}
\caption{Test macro-F1 per task. Bars are means over seeds; whiskers span the lowest and
highest seed.}
\label{fig:scores}
\end{figure}

\textit{Rail surface.} AST reaches 0.983 and 1.000 on its two seeds, above the forest
($0.958$), so the predictor uses AST. Over both seeds it misclassified only two of 78 wheel-burn
clips, both as normal track, and none of the others. Window length matters here: in the shared-model runs
(Section~\ref{sec:multitask}), 1\,s windows were too short to tell wheel burns from normal track,
whereas the separate AST uses 5\,s windows.

\textit{Corrugation at an unseen speed.} Both sensors detect corrugation on the 75\,km/h test run
after training on 40 and 50\,km/h. The axle-box forest is the most consistent model in the
study ($0.92 \pm 0.01$). It finds every corrugated test clip in all three seeds (recall 1.00),
at a recall of 0.89 on normal track. For the microphone, AST is ahead of the forest on average but
varies more between seeds: its recall on corrugated clips is 0.92 in one seed and 0.79 in the
other, with 0.93 on normal track in both. In both seeds its validation score was highest after the first epoch
and then fell until early stopping (Fig.~\ref{fig:curves}), so it overfits quickly on this small
set.

\textit{Thresholds matter.} With the default 0.5 threshold, the third microphone forest
predicted every test clip as normal (macro-F1 0.39); with the threshold tuned on the validation
run it scored 0.90. Every result here uses the tuned threshold.

\textit{Joints and switches.} This is the weakest task. The forest finds nearly all joint clips
(recall 0.96--0.98) but calls a third to half of the normal clips joints (normal recall
0.55--0.66; Fig.~\ref{fig:confusion}). AST trained on the same clips did no better in the
multi-task runs (Section~\ref{sec:multitask}). We think the labels are the limit: a normal clip
can contain an unlisted feature of the track, and the 6\,m labelling radius around each anomaly
is our choice, not a measured extent of the joint.

\begin{figure*}[t]
\centering
\fig{\textwidth}{confusion.pdf}
\caption{Test confusion matrices of the model the predictor uses for each task, summed over
seeds. Colour is the share of each true class; numbers are clip counts.}
\label{fig:confusion}
\end{figure*}

\begin{figure}[t]
\centering
\fig{\columnwidth}{training_curves.pdf}
\caption{Validation macro-F1 per epoch of the two AST models, one line per seed.}
\label{fig:curves}
\end{figure}

%% ------------------------------------------------------------------------------------------
\section{A Shared Model Did Not Help}
\label{sec:multitask}

Our first design was one shared AST encoder with a small output head per task, trained on all
four tasks at once. Each head compares a fault only with normal track from its own dataset, so
the model cannot score well by recognising which dataset a clip came from. Accelerometer signals
were treated as if sampled at 4\,kHz so that their content falls within AST's frequency range
while a clip stays longer than the window. Each batch picked a task, then a class, then a
recording.

Over four runs we changed the window length (1\,s to 2\,s), the accelerometer speed-up, the
training length, gave the microphone corrugation its own head, weighted the rail-surface task
twice, and kept the checkpoint whose weakest task scored best on validation rather than the
best mean. Only one run matched a separate model on any task (rail surface in run 2), and in that same run
microphone corrugation fell to 0.75 (Table~\ref{tab:mt}, Fig.~\ref{fig:mt}). The rail-surface task was the most fragile: in runs 3 and 4 the model all but
stopped predicting wheel burns, and macro-F1 fell below 0.6. When axle and microphone
corrugation shared one head, the microphone score dropped as the axle score rose. The joint task
stayed at 0.66--0.68 in every run, below the forest.

\begin{table}[t]
\centering
\caption{Test macro-F1 of the shared multi-task model in four runs, against the separate model
the predictor uses (mean over seeds).}
\label{tab:mt}
\begin{tabular}{lccccc}
\toprule
Task & Separate & Run 1 & Run 2 & Run 3 & Run 4 \\
\midrule
Rail surface       & \textbf{0.99} & 0.76 & \textbf{0.99} & 0.56 & 0.59 \\
Corrugation (axle) & \textbf{0.92} & 0.83 & 0.91 & 0.88 & 0.86 \\
Corrugation (mic)  & \textbf{0.89} & 0.85 & 0.75 & 0.82 & 0.88 \\
Joint / switch     & \textbf{0.81} & 0.67 & 0.68 & 0.66 & 0.68 \\
\bottomrule
\end{tabular}
\end{table}

\begin{figure*}[t]
\centering
\fig{\textwidth}{multitask_runs.pdf}
\caption{Test macro-F1 of each shared-model run (solid) against the separate model the
predictor uses (dashed).}
\label{fig:mt}
\end{figure*}

With datasets this small and this different, the tasks compete for the shared encoder rather
than help each other. Separate models are also simpler to retrain one at a time when new data
arrives for one fault, which is why we package them as one predictor instead.

%% ------------------------------------------------------------------------------------------
\section{Limitations}

Each dataset covers one line, one vehicle and one sensor setup, so none of these scores shows
that a model transfers to other track or other trains. The corrugation labels mark only the two
sections reported as severely corrugated; ``normal'' means not reported as corrugated rather than
inspected and clean. The joint labels come from surveyed positions and an unsupervised
detector, and joint scores are the weakest in the study. The AST models were trained with two
seeds rather than three because of limited GPU time. Finally, the four tasks come from different
recordings, so the predictor has not been evaluated on a single recording that contains several
kinds of fault.

%% ------------------------------------------------------------------------------------------
\section{Conclusion}

Public vibration and sound data is enough to build a working predictor for four rail-fault
tasks, reaching test macro-F1 of 0.99 on rail-surface faults, 0.92 and 0.89 for corrugation at an
unseen speed, and 0.81 on joints. Two simple steps did most of the work: cutting clips by
distance rather than time, and tuning the decision threshold on validation data. A shared
multi-task model did not match the separate models. The clearest next steps are better joint
labels, recordings from more lines and vehicles, and a test on a train carrying both a
microphone and accelerometers.

\section*{Data and Code Availability}
\label{sec:availability}
All datasets are public~\cite{shafique2021,upmcorrugation,ma2026railvivid}. Code to prepare the
clips, train every model, rebuild the tables and figures, and run the predictor is at
\url{https://github.com/Zcodeod-OG/Research-Projects} (folder \texttt{Rail-V/vibration}).

\bibliographystyle{IEEEtran}
\bibliography{refs}

\end{document}
